%% Chapter 1 — Introduction
%% OLDIA Thesis: Operational Lifecycle Determinism in Intelligent Architectures
%% Verified implementation facts as of 2026-06-09

\chapter{Introduction}
\label{chap:introduction}

\epigraph{%
  ``The question is not whether a machine can think, but whether we can \emph{prove} it thought
  correctly while it was doing so.''
}{--- paraphrase of the operational falsifiability problem}

Classical artificial intelligence produced architectures of genuine intellectual beauty.
MYCIN's certainty-factor algebra, Hearsay-II's blackboard, STRIPS's precondition--effect
formalism, and SOAR's preference-based decision cycle each resolved a hard theoretical
problem.
Yet every one of these systems shipped as a black box.
The 1976 MYCIN deployment had no audit trail.
Hearsay-II's knowledge-source activation record was discarded at runtime.
A correct diagnosis and an incorrect one looked identical from outside the process.
The epistemological gap between \emph{theoretical soundness} and \emph{operational
verifiability} was, for forty years, accepted as a fact of life.

This thesis argues that the gap is no longer necessary and that closing it changes the
economics of deploying certified reasoning at Fortune-5 scale.

%%─────────────────────────────────────────────────────────────────────────────
\section{The Epistemological Crisis in Deployed AI}
\label{sec:epistemological-crisis}
%%─────────────────────────────────────────────────────────────────────────────

\subsection{Theory without operational proof}

Shortliffe and Buchanan's certainty-factor formula~\cite{shortliffe1975mycin} is
mathematically well-defined.
Given two independent evidence items $e_1$ and $e_2$ and a hypothesis $h$, the combined
certainty factor obeys:
\[
  \mathrm{CF}(h, e_1 \wedge e_2) =
  \begin{cases}
    \mathrm{CF}_1 + \mathrm{CF}_2\,(1 - \mathrm{CF}_1) & \text{if both positive} \\
    \mathrm{CF}_1 + \mathrm{CF}_2\,(1 + \mathrm{CF}_2) & \text{if both negative} \\
    \dfrac{\mathrm{CF}_1 + \mathrm{CF}_2}{1 - \min(|\mathrm{CF}_1|, |\mathrm{CF}_2|)} & \text{otherwise}
  \end{cases}
\]
The formula is sound.
What it does not provide is any mechanism for an external auditor to verify, after the
fact, that the running system applied this formula to the correct inputs in the correct
order, that no rule was silently skipped, and that the final recommendation was derived
from a process consistent with the declared model.

This is not a criticism of Shortliffe and Buchanan.
Runtime process provability was not a design goal in 1976 because there was no
infrastructure to support it.
The insight of this thesis is that such infrastructure now exists --- in the form of
object-centric event logs, WebAssembly's deterministic execution model, and
cryptographic receipt chains --- and that retrofitting it onto classical architectures
reveals both their robustness and the depth of the original design decisions.

\subsection{The deployment gap in practice}

The gap between theory and deployment shows up in three ways that matter to practitioners.

\textbf{Silent deviation.}
A rule-based system that has been modified, patched, or ported may execute a different
process than its specification describes.
Without a conformance check, operators cannot distinguish a correct execution from a
deviant one.
Van der Aalst's work on process mining~\cite{vanderaalst2016process} formalised this
problem for business processes; this thesis applies the same discipline to AI reasoning
architectures.

\textbf{Unstable audit trails.}
Regulated industries (finance, pharmaceutical, critical infrastructure) require audit
evidence that a decision process conformed to a declared model.
Classical AI systems produce result logs, not process logs.
A result log proves that a recommendation was emitted; it does not prove that the
recommendation was reached by the declared reasoning process.

\textbf{Unverifiable latency claims.}
Microsecond-latency AI is increasingly demanded in high-frequency trading, real-time
threat detection, and autonomous systems.
A system that achieves low latency by skipping reasoning steps --- rather than by
executing them efficiently --- is operationally indistinguishable from one that does not,
unless the execution process is independently verified.

All three problems share a common structure: the \emph{output} is observable, but the
\emph{process} that produced it is not.
This thesis proposes and validates a framework that makes the process observable, falsifiable,
and cryptographically attested.

%%─────────────────────────────────────────────────────────────────────────────
\section{The Research Questions}
\label{sec:research-questions}
%%─────────────────────────────────────────────────────────────────────────────

Three research questions organise this investigation.

\begin{description}

\item[\textbf{RQ1}] \emph{Can classical AI reasoning be made operationally falsifiable ---
not just tested at development time, but continuously proven conforming at runtime?}

This question requires distinguishing two kinds of evidence.
A test suite proves that the implementation behaved correctly on a finite set of inputs
under controlled conditions.
Operational falsifiability means that, for every execution in production, an independent
artifact --- the OCEL event log and its associated BLAKE3 receipt --- can be presented to
an auditor who had no prior knowledge of the system and who can verify, from that artifact
alone, that the execution followed the declared lifecycle model.
Chapter~\ref{chap:ocel-provability} presents the formal model; Chapter~\ref{chap:validation}
presents evidence from 655 tests that the implementation meets the specification.

\item[\textbf{RQ2}] \emph{Does operational falsifiability change the deployment economics
for Fortune-5 enterprises?}

The hypothesis is that certified reasoning --- reasoning for which conformance evidence
is continuously produced and cryptographically attested --- commands a qualitatively
different position in the enterprise procurement conversation than uncertified reasoning
of equal accuracy.
Chapter~\ref{chap:deployment-patterns} develops a taxonomy of eight deployment patterns
that become commercially accessible when reasoning latency drops below 50 microseconds
and every execution produces a signed receipt.
The chapter argues that the combination of latency and provability unlocks markets that
are currently closed to AI: real-time regulatory compliance, intraday risk certification,
and autonomous systems with insurance-grade audit trails.

\item[\textbf{RQ3}] \emph{What becomes possible at 2--50 microsecond latency that the
original inventors could not have imagined?}

MYCIN was a batch consultation system.
Hearsay-II processed speech in near-real-time on 1970s hardware.
Neither system was designed for continuous, high-frequency production deployment.
When the same reasoning architectures execute in 2--50 microseconds inside a WebAssembly
sandbox, new application classes emerge: per-transaction fraud scoring with interpretable
rule traces, per-packet network policy reasoning with signed receipts, and per-event
process monitoring that emits OCEL conformance reports faster than the underlying business
process can advance.
Chapter~\ref{chap:deployment-patterns} catalogues these patterns and the conditions under
which they are achievable.

\end{description}

%%─────────────────────────────────────────────────────────────────────────────
\section{Contributions}
\label{sec:contributions}
%%─────────────────────────────────────────────────────────────────────────────

This thesis makes five original contributions.

\subsection{The wasm4pm cognition breed framework}

The central artifact of this thesis is a collection of thirteen validated, operationally
falsifiable implementations of classical AI architectures, compiled to WebAssembly and
exposed through a uniform \texttt{cognition\_run} interface.
The thirteen \emph{breeds} are: MYCIN (certainty-factor production rules), Hearsay-II
(blackboard with KSAR priority queue), SOAR (preference-algebra decision cycle with
chunking), CBR (case-based reasoning with full 4R cycle), Prolog (Robinson unification
with bounded SLD resolution), STRIPS (precondition--effect planning with soundness
verification), GPS (general problem solver with means--ends analysis), DENDRAL
(hypothesis generation with forbidden-candidate elimination), ELIZA (pattern-directed
dialogue), and four Autoinstinct breeds (vision, semantics, neurosis, learning).

Each breed is implemented from primary sources --- Shortliffe and Buchanan for MYCIN,
Newell for SOAR and GPS, Fikes and Nilsson for STRIPS, and so on --- and validated
against mathematical invariants that rule out stub implementations.
The Prolog breed, for example, must derive \texttt{grandparent(?0,?2)} from
\texttt{parent(?0,?1) :- ...} chains using genuine shared-variable unification; no
pattern-matching shortcut can satisfy this test.
The full test suite comprises 655 tests (480 Rust, 175 TypeScript), all passing as of
2026-06-09.

\subsection{The OCEL provability layer}

Every breed execution produces an OCEL 2.0~\cite{ghahfarokhi2021ocel} event log
capturing the reasoning trace as a sequence of typed events over named objects.
Thirteen lifecycle models --- one per breed --- define the lawful sequences of event types
for that architecture.
A DFA-based conformance checker (\texttt{validate\_ocel\_alignment}) verifies, at the
execution choke point in \texttt{wasm.rs}, that every emitted trace conforms to the
declared model before the result is returned to the caller.
Non-conforming executions are refused, not logged for later review.

The OCEL layer is implemented in \texttt{crates/wasm4pm-cognition/src/ocel.rs} (632
lines).
Logical-step attributes rather than wall-clock timestamps ensure that conformance
evidence is independent of hardware timing and replayable in any environment.

\subsection{The BLAKE3 receipt chain}

Every successful execution emits a cryptographic receipt.
The \texttt{output\_hash} is \texttt{blake3(serde\_json(BreedOutput))}, where
\texttt{BreedOutput} contains the OCEL log.
The receipt therefore attests to \emph{process}, not just result: an auditor who holds
the receipt can verify that the output was produced by a specific reasoning trace, not
merely that some output was produced.
The \texttt{run\_id} is \texttt{blake3(breed~||~output\_hash)}, making receipts
cross-breed comparable.
The \texttt{replay\_pointer} (the first 16 hex characters of \texttt{output\_hash})
provides a compact, stable identifier suitable for embedding in external audit systems.

\subsection{The unfakeable oracle test suite}

The 655-test suite is designed so that no stub implementation --- no implementation that
returns hardcoded outputs without executing the declared algorithm --- can pass.
This property, which we call \emph{oracle unfakeability}, is achieved by testing
mathematical invariants rather than specific output values.
MYCIN tests the strict boundary at CF = 0.2 (rejected) and CF = 0.201 (accepted), which
requires the continuous formula to be evaluated.
CBR tests that identical query and case produce Jaccard similarity of exactly 1.0.
SOAR tests that circular worse-than preferences produce deterministic termination, which
requires the full preference algebra to execute.
AutoinstinctLearning tests that plan-step distance to goal is non-increasing, which
requires genuine planning.

\subsection{A taxonomy of Fortune-5 deployment patterns}

Chapter~\ref{chap:deployment-patterns} presents eight deployment patterns enabled by the
combination of microsecond latency and operational provability.
The patterns are grounded in the verified timing characteristics of the breed framework
(2--50 microseconds per execution on 2026 hardware) and in the structural requirements
of the OCEL provability and BLAKE3 receipt layers.
Each pattern is described in terms of the reasoning architecture it employs, the
conformance evidence it produces, and the commercial context in which it is valuable.

%%─────────────────────────────────────────────────────────────────────────────
\section{Document Structure}
\label{sec:document-structure}
%%─────────────────────────────────────────────────────────────────────────────

The remainder of this thesis is organised as follows.

Chapter~\ref{chap:background} reviews the classical AI architectures that form the
substrate of the breed framework, the process mining literature that motivates the OCEL
provability layer, and prior work on certified AI and runtime verification.

Chapter~\ref{chap:breed-framework} describes the thirteen breed implementations in
detail: their theoretical foundations, their Rust/WASM implementations, and the
design decisions that make them simultaneously faithful to the originals and amenable
to operational falsifiability.

Chapter~\ref{chap:ocel-provability} presents the OCEL provability layer: the
object-centric event log schema, the thirteen lifecycle models, the DFA conformance
checker, and the temporal conformance enforcement mechanism.

Chapter~\ref{chap:receipt-chain} presents the BLAKE3 receipt chain: its construction,
its relationship to the OCEL log, and its properties as an audit instrument.

Chapter~\ref{chap:validation} presents the validation evidence: the 655-test suite,
the oracle unfakeability argument, and the timing measurements that establish the
microsecond-latency claims.

Chapter~\ref{chap:deployment-patterns} presents the taxonomy of Fortune-5 deployment
patterns and argues that operational falsifiability is a necessary condition for
their commercial viability.

Chapter~\ref{chap:conclusion} summarises the contributions, states the limitations of
the current implementation, and identifies directions for future work.

%%─────────────────────────────────────────────────────────────────────────────
\section{Terminology}
\label{sec:terminology}
%%─────────────────────────────────────────────────────────────────────────────

The following terms are used with precise meanings throughout this thesis.

\begin{description}

\item[\textbf{Breed.}]
A \emph{breed} is a validated, operationally falsifiable implementation of a classical AI
reasoning architecture, compiled to WebAssembly and exposed through the uniform
\texttt{cognition\_run} interface.
The term \emph{breed} is preferred over \emph{algorithm} or \emph{agent} because each
implementation encompasses both the reasoning procedure and the lifecycle model against
which its execution is continuously verified.
Thirteen breeds are implemented in this thesis.

\item[\textbf{OCEL.}]
\emph{OCEL} (Object-Centric Event Log)~\cite{ghahfarokhi2021ocel} is a log format in
which events are related to multiple objects of multiple types, rather than to a single
case identifier.
OCEL 2.0 extends the original format with typed attributes on event--object relationships.
This thesis uses OCEL 2.0 as the provability substrate: each breed execution emits an
OCEL log in which the reasoning objects (rules, knowledge sources, goals, cases) are
first-class entities with traceable lifecycles.

\item[\textbf{Lifecycle model.}]
A \emph{lifecycle model} is a finite automaton that defines the lawful sequences of event
types for a given breed.
For example, the MYCIN lifecycle model requires that every execution begin with
\texttt{rule\_evaluate} events, followed by \texttt{cf\_combine} events, and conclude
with a \texttt{threshold\_check} event.
A trace that omits the \texttt{threshold\_check} or emits events in the wrong order is
non-conforming and is refused at the execution choke point.

\item[\textbf{Conformance.}]
\emph{Conformance} is the binary property of an OCEL trace with respect to a lifecycle
model: a trace either conforms or it does not.
This thesis does not use graded conformance metrics (fitness, precision) in the
execution-time path; those metrics appear in the validation chapter where the test suite
is analysed.
At the choke point, conformance is a gate: non-conforming traces are refused.

\item[\textbf{Receipt chain.}]
The \emph{receipt chain} is the sequence of BLAKE3-hashed receipts produced by successive
executions of a breed.
Each receipt's \texttt{output\_hash} covers the OCEL log, making the chain a
tamper-evident record of the reasoning process, not merely of the results.
The term \emph{chain} reflects the structural dependency: \texttt{run\_id} is derived
from \texttt{output\_hash}, which is derived from the \texttt{BreedOutput}, which
contains the OCEL log.

\item[\textbf{Unfakeable oracle.}]
An \emph{unfakeable oracle} is a test assertion that can only be satisfied by an
implementation that executes the declared algorithm.
It is distinct from a value-equality assertion (which tests a specific output) and from
a property-based assertion (which tests a statistical distribution of outputs).
An unfakeable oracle tests a mathematical invariant that is \emph{structurally impossible}
to satisfy without executing the algorithm: for example, the Jaccard identity invariant
for CBR, or the shared-variable unification requirement for Prolog.

\item[\textbf{KSAR.}]
A \emph{KSAR} (Knowledge Source Activation Record) is the scheduling primitive of the
Hearsay-II blackboard architecture~\cite{erman1980hearsay}.
When a knowledge source detects a triggering condition on the blackboard, it creates a
KSAR recording its identity, the triggering evidence, and an estimated priority
(\texttt{rating = certainty * trigger\_cf}).
The blackboard scheduler selects the highest-rated pending KSAR for execution.
In the wasm4pm implementation, KSARs are managed in a priority queue using
\texttt{f32::total\_cmp} for deterministic ordering, and stale KSARs are re-rated when
the blackboard state changes.

\end{description}

%%─────────────────────────────────────────────────────────────────────────────
%% End of Chapter 1
%%─────────────────────────────────────────────────────────────────────────────
